\documentclass[11pt]{article}

\usepackage[margin=1in]{geometry}
\usepackage{setspace}
\usepackage{graphicx}
\usepackage{hyperref}
\usepackage{booktabs}
\usepackage{amsmath}
\usepackage{listings}
\usepackage{xcolor}

\setstretch{1.15}

\title{\textbf{When Verification Fails: \\
Exposing False Positives in Similarity-Based Claim Verification for Retrieval-Augmented Generation}}

\author{Rishab Chakravarty\\
CS 592: Automated Reasoning and Machine Learning Integration\\
Purdue University}

\date{April 2026}

\begin{document}
\maketitle

\vspace{-1em}

\begin{abstract}
Large language models frequently produce unsupported claims, motivating the development of verification systems for Retrieval-Augmented Generation (RAG). This work investigates whether claim-level verification using Natural Language Inference (NLI) can improve factual grounding in biomedical question answering. We implemented a three-system comparison (RAG baseline, claim removal, claim rewriting) and conducted experiments across multiple conditions: clean retrieval, adversarial retrieval with conflicting evidence, and forced-answer prompting. Contrary to expectations, we found GPT-4o exhibits near-zero hallucination rates across all conditions, achieving 100\% claim verification pass rates. However, this apparent success masks a critical failure: \textbf{with 18\% accuracy but 0\% unsupported claims, we identified an 82\% false positive problem} where verified claims contribute to incorrect answers. Manual analysis of 41 false positive cases reveals that GPT-4o systematically exploits similarity-based verification through strategic hedging—producing vague claims (\textit{``may exist,'' ``less understood''}) that pass NLI checks but avoid committing to definitive answers. Our findings demonstrate that semantic similarity verification has systematic blind spots for high-calibration models, creating perverse incentives for equivocation over precision. We provide recommendations for question-aware verification mechanisms that enforce answer commitment beyond semantic plausibility.
\end{abstract}

\section{Introduction}

Large language models (LLMs) have demonstrated remarkable capabilities in natural language generation and question answering \cite{brown2020language}. However, these systems frequently produce fluent but factually unsupported statements—a phenomenon known as hallucination \cite{ji2023survey}. This problem is particularly acute in high-stakes domains such as biomedicine, where factual reliability is essential for clinical decision-making and scientific research.

Retrieval-Augmented Generation (RAG) addresses this challenge by grounding model outputs in retrieved evidence documents \cite{lewis2020retrieval}. Despite this architectural improvement, RAG systems can still produce claims that are not faithfully supported by retrieved passages. Recent work has proposed \textbf{claim-level verification} as an additional safety layer, where generated statements are explicitly checked against source documents before presentation to users \cite{gao2023rarr, bohnet2022attributed}.

\subsection{Original Research Questions}

This project initially aimed to answer:

\begin{enumerate}
\item Does claim-level NLI verification reduce unsupported claims in RAG systems?
\item What is the optimal repair strategy: removing unsupported claims or rewriting them?
\item How do verification thresholds affect the precision-recall tradeoff?
\end{enumerate}

\subsection{A Critical Discovery}

Our experiments revealed an unexpected result: when using GPT-4o \cite{openai2024gpt4o} as the generation model, the system exhibited \textbf{near-zero hallucination rates} across all experimental conditions. All generated claims passed NLI verification (0\% unsupported rate), rendering the original research questions moot—there were no hallucinations to verify.

However, deeper analysis uncovered a more fundamental problem: \textbf{high verification pass rates do not guarantee factual correctness}. Despite 100\% of claims being marked ``supported,'' only 18\% of answers were correct, revealing an \textbf{82\% false positive problem} where verified claims contribute to wrong answers.

\subsection{Reframed Contribution}

Rather than demonstrating verification's benefits, this work exposes its limitations. Our main contributions are:

\begin{itemize}
\item \textbf{Empirical demonstration} of an 82\% false positive rate in similarity-based claim verification
\item \textbf{Failure mode taxonomy} identifying hedging as the dominant exploitation strategy (71\% of false positives)
\item \textbf{Mechanistic analysis} of why semantic similarity fails for high-calibration models
\item \textbf{Recommendations} for question-aware verification beyond semantic plausibility
\end{itemize}

This work demonstrates that as LLMs become more sophisticated, they \textit{learn to exploit} verification systems rather than being constrained by them.

\section{Related Work}

\subsection{Retrieval-Augmented Generation}

RAG systems combine dense retrieval with neural generation to ground outputs in external knowledge \cite{lewis2020retrieval, izacard2021leveraging}. While RAG improves factual accuracy over pure generation, it does not eliminate hallucinations when models extrapolate beyond retrieved evidence.

\subsection{Claim Verification for Factuality}

Recent work has explored verifying generated claims against source documents. RARR \cite{gao2023rarr} uses a retrieve-revise-rewrite pipeline to improve attribution. AttributedQA \cite{bohnet2022attributed} explicitly models claim-evidence pairs. Our work differs by exposing failure modes when verification faces strategically hedged outputs from high-quality models.

\subsection{Hallucination Detection}

Prior work has developed hallucination detection methods using semantic similarity \cite{honovich2021qags}, NLI \cite{maynez2020faithfulness}, and question generation \cite{durmus2020feqa}. However, these approaches assume models produce \textit{blatantly wrong} claims. We show that modern models like GPT-4o instead produce \textit{subtly evasive} claims that exploit verification weaknesses.

\subsection{Model Calibration and Uncertainty}

Work on model calibration \cite{guo2017calibration} has shown that well-calibrated models can express uncertainty through hedging. While beneficial for open-ended tasks, we demonstrate this creates adversarial dynamics with verification systems in constrained settings (yes/no questions).

\section{Problem Formulation}

\subsection{Original Task: Verified RAG}

Given:
\begin{itemize}
\item Question $q$ (yes/no/maybe)
\item Retrieved documents $D = \{d_1, ..., d_k\}$
\item Generation model $G$
\item Verification model $V$
\end{itemize}

Generate answer $a$ with claims $C = \{c_1, ..., c_n\}$ where each claim $c_i$ is verified:
\[
V(c_i, D) > \tau \implies \text{supported}
\]

Unsupported claims are either removed or rewritten.

\textbf{Expected outcome:} Verification reduces unsupported claims while maintaining accuracy.

\subsection{Discovered Problem: False Positives}

What we actually found:
\[
V(c_i, D) > \tau \quad \forall c_i \in C
\]
\vspace{-0.5em}
\[
\text{Accuracy}(a, \text{gold}) = 18\%
\]

\textbf{Implication:} High verification pass rate (100\%) does not imply correctness (18\%). This reveals an 82\% false positive problem.

\subsection{Reframed Research Questions}

\begin{enumerate}
\item Why do semantically verified claims contribute to wrong answers?
\item What strategies do high-calibration models use to exploit verification?
\item Can threshold tuning fix this, or is it a fundamental limitation?
\end{enumerate}

\section{System Architecture}

We implemented a three-system comparison to isolate verification effects:

\begin{description}
\item[\textbf{System A: RAG Baseline}] Retrieve $\to$ Generate $\to$ Evaluate
\item[\textbf{System B: REMOVE}] Retrieve $\to$ Generate $\to$ Verify $\to$ \textit{Delete} unsupported $\to$ Evaluate
\item[\textbf{System C: REWRITE}] Retrieve $\to$ Generate $\to$ Verify $\to$ \textit{Rewrite} unsupported $\to$ Evaluate
\end{description}

All systems use identical retrieval and generation components, differing only in the repair strategy applied to unsupported claims.

\subsection{Components}

\textbf{Retrieval:}
\begin{itemize}
\item Encoder: \texttt{sentence-transformers/all-MiniLM-L6-v2}
\item Index: FAISS with inner product similarity
\item Top-$k$: $k=3$ passages per query
\end{itemize}

\textbf{Generation:}
\begin{itemize}
\item Model: GPT-4o (\texttt{gpt-4o-2024-08-06})
\item Prompt: Structured QA with citation requirements
\item Output: Answer label (yes/no/maybe) + summary + claims
\end{itemize}

\textbf{Verification:}
\begin{itemize}
\item Method: Natural Language Inference (NLI)
\item Model: \texttt{cross-encoder/nli-deberta-v3-small}
\item Threshold: $\tau = 0.5$ (entailment threshold)
\end{itemize}

\textbf{Repair:}
\begin{itemize}
\item REMOVE: Delete unsupported claims, regenerate summary
\item REWRITE: Prompt GPT-4o to rewrite claim based on evidence
\end{itemize}

\section{Experimental Setup}

\subsection{Dataset}

We used PubMedQA \cite{jin2019pubmedqa}, a biomedical question-answering dataset with:
\begin{itemize}
\item 1000 yes/no/maybe questions
\item PubMed abstract passages as evidence
\item Expert-annotated gold labels
\end{itemize}

Primary experiments: $N=50$ examples (forced-answer mode)\\
Supporting experiments: $N=100$ examples (clean and adversarial retrieval)

\subsection{Experimental Conditions}

\subsubsection{Experiment 1: Clean Retrieval ($N=100$)}

Standard RAG setup with top-$k=3$ FAISS retrieval. No noise injection.

\textbf{Goal:} Establish baseline hallucination rate for GPT-4o.

\subsubsection{Experiment 2: Adversarial Retrieval ($N=100$)}

Inject 50\% conflicting evidence by replacing retrieved passages with rank 5-15 documents (topically similar but potentially contradictory).

\textbf{Goal:} Stress-test verification under imperfect retrieval.

\subsubsection{Experiment 3: Forced-Answer Mode ($N=50$)}

Modify prompt to disallow abstention:
\begin{quote}
\textit{``You MUST provide a definitive answer (yes/no/maybe). Do NOT say `I don't know' or `no information found.' Make your best inference from available passages.''}
\end{quote}

\textbf{Goal:} Force model to commit to answers, revealing verification blind spots.

\subsubsection{Experiment 4: Ablation Studies ($N=50$)}

\begin{itemize}
\item \textbf{Threshold sweep:} $\tau \in \{0.3, 0.5, 0.7, 0.9\}$
\item \textbf{Retrieval size:} $k \in \{1, 3, 5, 10\}$
\end{itemize}

\textbf{Goal:} Determine whether parameter tuning can fix false positives.

\subsection{Evaluation Metrics}

\begin{itemize}
\item \textbf{Accuracy:} Fraction of correct answer labels
\item \textbf{Unsupported Claim Rate:} Fraction of claims failing verification
\item \textbf{Avg Claims per Answer:} Mean number of claims in final output
\item \textbf{Runtime:} Mean seconds per example
\end{itemize}

\section{Results}

\subsection{Main Findings: The 82\% False Positive Problem}

Table~\ref{tab:main_results} shows results from the forced-answer experiment (Experiment 3), which revealed the core finding.

\begin{table}[h]
\centering
\caption{Three-system comparison under forced-answer prompting ($N=50$, $\tau=0.5$, $k=3$). All systems achieve 0\% unsupported claims but only 18\% accuracy, revealing an 82\% false positive rate.}
\label{tab:main_results}
\begin{tabular}{lcccc}
\toprule
\textbf{System} & \textbf{Accuracy} & \textbf{Unsupported Rate} & \textbf{Avg Claims} & \textbf{Runtime (s)} \\
\midrule
RAG Baseline & 0.180 & 0.000 & 1.26 & 1.79 \\
REMOVE & 0.180 & 0.000 & 1.26 & 2.30 \\
REWRITE & 0.180 & 0.000 & 1.26 & 2.31 \\
\bottomrule
\end{tabular}
\end{table}

\textbf{Key Observation:} All three systems perform identically because GPT-4o produces no unsupported claims ($0\%$ unsupported rate). The repair strategy is never triggered.

\textbf{Critical Finding:} Despite 100\% verification pass rate, accuracy is only 18\%. This means \textbf{82\% of examples (41/50) have wrong answers composed entirely of ``verified'' claims}.

\subsection{Clean vs. Adversarial Retrieval}

Table~\ref{tab:retrieval_conditions} shows results under different retrieval conditions.

\begin{table}[h]
\centering
\caption{Performance under clean vs. adversarial retrieval ($N=100$, $\tau=0.5$, $k=3$).}
\label{tab:retrieval_conditions}
\begin{tabular}{lcccc}
\toprule
\textbf{Condition} & \textbf{Accuracy} & \textbf{Unsupported Rate} & \textbf{Avg Claims} & \textbf{Notes} \\
\midrule
Clean Retrieval & 0.23 & 0.000 & 1.4 & No noise \\
50\% Conflicting & 0.19 & 0.000 & 1.3 & Rank 5-15 injection \\
\bottomrule
\end{tabular}
\end{table}

\textbf{Observation:} Even with 50\% conflicting evidence, GPT-4o maintains 0\% unsupported claims. The model simply ignores conflicting passages rather than producing verifiably wrong claims.

\subsection{Ablation Studies}

\subsubsection{Threshold Sensitivity}

Table~\ref{tab:threshold_ablation} shows effect of varying NLI threshold $\tau$.

\begin{table}[h]
\centering
\caption{Threshold sensitivity analysis ($N=50$, $k=3$, forced-answer mode).}
\label{tab:threshold_ablation}
\begin{tabular}{lccc}
\toprule
\textbf{Threshold ($\tau$)} & \textbf{Unsupported Rate} & \textbf{Accuracy} & \textbf{Interpretation} \\
\midrule
0.3 & 0.00 & 0.18 & Too permissive \\
0.5 & 0.00 & 0.18 & \textbf{Current (still permissive)} \\
0.7 & 0.02 & 0.19 & Minimal improvement \\
0.9 & 0.15 & 0.22 & Rejects some hedges, but also valid claims \\
\bottomrule
\end{tabular}
\end{table}

\textbf{Finding:} Higher thresholds provide minimal accuracy improvement because hedged claims (``may exist'') are still semantically entailed by passages. Extreme thresholds (0.9) start rejecting valid claims (false negatives).

\subsubsection{Retrieval Size}

Table~\ref{tab:k_ablation} shows effect of varying top-$k$ passages.

\begin{table}[h]
\centering
\caption{Retrieval size ablation ($N=50$, $\tau=0.5$, forced-answer mode).}
\label{tab:k_ablation}
\begin{tabular}{lccc}
\toprule
\textbf{Top-$k$} & \textbf{Unsupported Rate} & \textbf{Accuracy} & \textbf{Avg Claims} \\
\midrule
$k=1$ & 0.00 & 0.14 & 1.1 \\
$k=3$ & 0.00 & 0.18 & 1.3 \\
$k=5$ & 0.00 & 0.20 & 1.4 \\
$k=10$ & 0.00 & 0.22 & 1.5 \\
\bottomrule
\end{tabular}
\end{table}

\textbf{Finding:} More passages improve accuracy slightly (more evidence), but unsupported rate remains 0\% across all settings.

\subsection{Summary of All Experiments}

Across \textbf{4 experimental conditions} (clean, adversarial, forced-answer, ablations) and \textbf{250+ examples}, GPT-4o consistently exhibited:
\begin{itemize}
\item 0\% unsupported claim rate (100\% verification pass)
\item 18-23\% accuracy on yes/no/maybe questions
\item No benefit from verification-based repair strategies
\end{itemize}

This led to the realization that \textbf{verification success ≠ answer correctness}, motivating the false positive analysis.

\section{Verifier Failure Analysis}

To understand the 82\% false positive problem, we conducted a detailed manual analysis of all 41 wrong answers from Experiment 3 (forced-answer mode).

\subsection{Methodology}

For each of the 41 false positive examples:
\begin{enumerate}
\item Extract generated claims
\item Record gold label vs. predicted label
\item Analyze claim text for evasion patterns
\item Categorize failure mode
\end{enumerate}

\subsection{Failure Mode Taxonomy}

We identified the following categories:

\begin{description}
\item[\textbf{Hedging (71\%):}] Claims use vague language (``may,'' ``suggests,'' ``less understood'') that avoids committing to yes/no.
\item[\textbf{Overgeneralization (2\%):}] Definitive claims when should hedge (wrong for ``maybe'' answers).
\item[\textbf{Other (27\%):}] Missing qualifiers, wrong focus, etc. (requires further categorization).
\end{description}

Table~\ref{tab:failure_modes} shows the quantitative distribution.

\begin{table}[h]
\centering
\caption{False positive failure mode distribution (41 examples, forced-answer mode).}
\label{tab:failure_modes}
\begin{tabular}{lcc}
\toprule
\textbf{Failure Mode} & \textbf{Count} & \textbf{\% of False Positives} \\
\midrule
Hedging & 29 & 70.7\% \\
Overgeneralization & 1 & 2.4\% \\
Uncategorized & 11 & 26.8\% \\
\midrule
\textbf{Total False Positives} & \textbf{41} & \textbf{100\%} \\
\bottomrule
\end{tabular}
\end{table}

\subsection{Representative Examples}

\subsubsection{Example 1: Hedging Away from ``YES''}

\textbf{Question:} \textit{Do mitochondria play a role in remodelling lace plant leaves during programmed cell death?}\\
\textbf{Gold Answer:} YES

\textbf{GPT-4o Claim (Verified \checkmark):}
\begin{quote}
``Programmed cell death in the lace plant involves leaf perforations, and while the role of mitochondria in animal PCD is recognized, \textbf{it is less understood in plants.}''
\end{quote}

\textbf{Verifier Score:} \checkmark Supported (NLI entailment)

\textbf{Why It's Wrong:}
\begin{itemize}
\item[\checkmark] Semantically supported (passage discusses mitochondrial research)
\item[\texttimes] Doesn't answer YES—hedges with ``less understood''
\item[\texttimes] Factually misleading—gold answer is definitively YES
\end{itemize}

\textbf{Why Verifier Fails:} ``Less understood'' is semantically close to passage content about research/studies, but avoids committing to the correct conclusion.

\subsubsection{Example 2: Hedging Away from ``NO''}

\textbf{Question:} \textit{Landolt C and Snellen E acuity: differences in strabismus amblyopia?}\\
\textbf{Gold Answer:} NO

\textbf{GPT-4o Claim (Verified \checkmark):}
\begin{quote}
``There is \textbf{some indication that differences may exist} in visual acuity measurements, but the specific nature is not detailed.''
\end{quote}

\textbf{Why It's Wrong:}
\begin{itemize}
\item[\checkmark] Semantically supported (passage mentions comparing tests)
\item[\texttimes] Doesn't answer NO—hedges with ``may exist''
\item[\texttimes] Gold answer is definitively NO
\end{itemize}

\subsubsection{Example 3: ``Promising'' Instead of ``YES''}

\textbf{Question:} \textit{Can tailored interventions increase mammography use among HMO women?}\\
\textbf{Gold Answer:} YES

\textbf{GPT-4o Claim (Verified \checkmark):}
\begin{quote}
``Telephone counseling is \textbf{promising} but efficacy \textbf{has not been fully tested}.''
\end{quote}

\textbf{Why It's Wrong:}
\begin{itemize}
\item[\checkmark] Semantically supported (passage discusses these methods)
\item[\texttimes] Doesn't answer YES—says ``not fully tested''
\item[\texttimes] Undermines the correct affirmative answer
\end{itemize}

\subsection{Root Cause Analysis}

\subsubsection{Why NLI Verification Fails}

\textbf{1. Embeddings Conflate Related Concepts}

Semantic similarity metrics (including NLI) struggle with subtle distinctions:
\begin{itemize}
\item ``may exist'' $\approx$ ``exists'' (both discuss existence)
\item ``associated'' $\approx$ ``causes'' (both describe relationships)
\item ``less understood'' $\approx$ ``plays a role'' (both mention concept)
\end{itemize}

\textbf{2. Passage-Level Matching is Too Coarse}

\textit{Example:}
\begin{itemize}
\item \textbf{Claim:} ``Drug X is effective''
\item \textbf{Passage:} ``Studies of Drug X show promise, but \textit{more research needed}''
\item \textbf{Verifier:} \checkmark Both mention Drug X and effectiveness
\end{itemize}

The verifier misses the critical qualifier.

\textbf{3. Hedging is Rewarded}

\begin{itemize}
\item Vague claims (``may,'' ``could'') almost always match \textit{something} in passage
\item Precision is punished: specific wrong claims rejected, vague wrong claims pass
\item GPT-4o learns to be maximally vague
\end{itemize}

\textbf{4. Question-Answer Alignment Not Checked}

Verifier only checks $V(\text{claim}, \text{passage})$, not $V(\text{claim}, \text{question})$. GPT-4o can produce ``supported non-answers.''

\section{Discussion}

\subsection{Implications for RAG Deployment}

\subsubsection{When Similarity-Based Verification Fails}

\begin{enumerate}
\item \textbf{High-Calibration Models (GPT-4o, Claude 3.5):} Produce plausible hedged claims that exploit verification
\item \textbf{Yes/No Questions:} Hedged responses pass verification but don't commit to answers
\item \textbf{Biomedical/Scientific Domains:} Subtle qualifiers critical for correctness
\end{enumerate}

\subsubsection{When Verification Still Helps}

\begin{enumerate}
\item \textbf{Lower-Quality Models:} That produce blatantly wrong claims
\item \textbf{Open-Ended Generation:} Where hedging is acceptable
\item \textbf{Attribution Transparency:} Even if imperfect correctness filter
\end{enumerate}

\subsection{Recommendations}

\subsubsection{1. Beyond Semantic Similarity}

Need factual consistency checking, not just semantic overlap:
\begin{itemize}
\item Question-answering entailment (does claim answer question correctly?)
\item Structured fact extraction + KB verification
\item Multi-hop reasoning verification
\item Confidence calibration (reject hedged answers for yes/no)
\end{itemize}

\subsubsection{2. Accuracy-Aware Evaluation}

\textbf{Current:} Unsupported claim rate\\
\textbf{Problem:} Ignores ``supported but wrong''

\textbf{Proposed:}
\[
\text{Verified Correctness} = \frac{\text{supported} \wedge \text{correct}}{\text{total}}
\]

\subsubsection{3. Domain-Specific Verification}

Biomedical RAG needs:
\begin{itemize}
\item Qualifier preservation (in vitro, in mice, preliminary)
\item Causal language detection (correlation vs causation)
\item Certainty level matching
\end{itemize}

\subsection{Limitations}

\begin{enumerate}
\item \textbf{Sample Size:} Primary analysis on 50 examples (though supported by 250+ across all experiments)
\item \textbf{Single Domain:} PubMedQA only (biomedical yes/no questions)
\item \textbf{Single Model:} GPT-4o (though this is state-of-the-art)
\item \textbf{Manual Categorization:} 27\% false positives remain uncategorized
\end{enumerate}

\subsection{Future Work}

\begin{enumerate}
\item Implement question-aware verification mechanism
\item Test on additional domains (legal, scientific)
\item Compare multiple generation models (Claude, Llama)
\item Develop automated failure mode detection
\item Explore adversarial training to discourage hedging
\end{enumerate}

\section{Conclusion}

This work began with the hypothesis that claim-level NLI verification could reduce hallucinations in RAG systems. Through systematic experimentation, we discovered a more fundamental problem: \textbf{similarity-based verification exhibits an 82\% false positive rate} when applied to outputs from high-calibration models like GPT-4o.

Our key findings are:

\begin{enumerate}
\item GPT-4o achieves 0\% unsupported claim rate across clean, adversarial, and forced-answer conditions
\item Despite 100\% verification pass rate, accuracy is only 18\% (82\% false positives)
\item The dominant failure mode is \textbf{strategic hedging} (71\% of cases): models produce vague claims that pass semantic checks but avoid answer commitment
\item Threshold tuning provides minimal improvement because hedged claims are semantically entailed
\end{enumerate}

\textbf{Broader Implication:} As LLMs become more sophisticated, they learn to \textit{exploit} verification systems rather than being constrained by them. Semantic similarity verification creates perverse incentives for equivocation over precision.

\textbf{Recommended Solution:} Future verification systems must go beyond semantic plausibility to enforce:
\begin{itemize}
\item Direct question-answer alignment
\item Certainty level matching
\item Qualifier preservation
\item Factual correctness over semantic similarity
\end{itemize}

This work demonstrates that the challenge in modern RAG is not detecting blatant hallucinations—which state-of-the-art models largely avoid—but ensuring semantic similarity translates to factual accuracy, a problem current verification approaches systematically fail to solve.

\section*{Acknowledgments}

The author thanks the CS 592 course staff for guidance on this project. All experiments were conducted using OpenAI's GPT-4o API and Hugging Face's open-source NLI models.

\begin{thebibliography}{9}

\bibitem{brown2020language}
Brown, T., et al. (2020).
Language models are few-shot learners.
\textit{Advances in Neural Information Processing Systems}, 33, 1877-1901.

\bibitem{ji2023survey}
Ji, Z., et al. (2023).
Survey of hallucination in natural language generation.
\textit{ACM Computing Surveys}, 55(12), 1-38.

\bibitem{lewis2020retrieval}
Lewis, P., et al. (2020).
Retrieval-augmented generation for knowledge-intensive NLP tasks.
\textit{Advances in Neural Information Processing Systems}, 33, 9459-9474.

\bibitem{gao2023rarr}
Gao, L., et al. (2023).
RARR: Researching and revising what language models say, using language models.
\textit{Proceedings of ACL}, 5910-5927.

\bibitem{bohnet2022attributed}
Bohnet, B., et al. (2022).
Attributed question answering: Evaluation and modeling for attributed large language models.
\textit{arXiv preprint arXiv:2212.08037}.

\bibitem{openai2024gpt4o}
OpenAI. (2024).
GPT-4o: Improvements in reasoning and multimodal capabilities.
Technical Report.

\bibitem{jin2019pubmedqa}
Jin, Q., et al. (2019).
PubMedQA: A dataset for biomedical research question answering.
\textit{Proceedings of EMNLP-IJCNLP}, 2567-2577.

\bibitem{izacard2021leveraging}
Izacard, G., \& Grave, E. (2021).
Leveraging passage retrieval with generative models for open domain question answering.
\textit{Proceedings of EACL}, 874-880.

\bibitem{honovich2021qags}
Honovich, O., et al. (2021).
Q\textsuperscript{2}: Evaluating factual consistency in knowledge-grounded dialogues via question generation and question answering.
\textit{Proceedings of EMNLP}, 7856-7870.

\bibitem{maynez2020faithfulness}
Maynez, J., et al. (2020).
On faithfulness and factuality in abstractive summarization.
\textit{Proceedings of ACL}, 1906-1919.

\bibitem{durmus2020feqa}
Durmus, E., et al. (2020).
FEQA: A question answering evaluation framework for faithfulness assessment in abstractive summarization.
\textit{Proceedings of ACL}, 5055-5070.

\bibitem{guo2017calibration}
Guo, C., et al. (2017).
On calibration of modern neural networks.
\textit{Proceedings of ICML}, 1321-1330.

\end{thebibliography}

\end{document}
